% Section 7.5, from lectures/L07/appendix/d_exercises.md.
%
% One correction to the source: exercise 10 a) asks for the cycles of exercise 4's program from
% start, and its solution counts the rjmp main at address 0, which the listing in exercise 4 does
% not show. The exercise now says that the program sits in the template, here and in the source.

\section{Övningar}
\label{c7:sec:exercises}\label{c7:app:d}

\begin{aside}
\textbf{Så kontrollerar du ditt arbete.} Övningarna märkta \textbf{Program} och \textbf{Kontroll}
görs i Microchip Studio: skapa ett projekt enligt bilaga~\ref{app:studio}, skriv programmet, bygg
med \textbf{F7} och stega med \textbf{F11}. Förutsäg alltid resultatet innan du stegar
(\secref{c7:c3}).

Övriga övningar görs med papper och penna. Lösningsförslag finns i bilaga~\ref{app:answers}. Gör
övningen innan du läser lösningen: det du nästan svarade lär dig mer än det rätta svaret gör.
\end{aside}

Varje övning är märkt med sin sort. Det finns en \textbf{Kontroll} i det här kapitlet,
övning~\ref{c7:ex:7}.

\exercise{Arbetsbänken}{Förståelse}
\label{c7:ex:1}
\begin{parts}
  \item Hur många register har ATmega328P, och hur många bitar rymmer vart och ett?
  \item Vilka register kan laddas med \code{ldi}?
  \item Förklara, med hjälp av hur instruktionen \code{ldi} är kodad, varför de övriga registren
    inte kan laddas med \code{ldi}. Svaret ska innehålla ett antal bitar.
  \item Du vill ha talet 7 i \code{r5}. Skriv två instruktioner som gör det.
\end{parts}

\exercise{Tre minnen}{Förståelse}
\label{c7:ex:2}
\begin{parts}
  \item I vilket minne ligger programmet?
  \item I vilket minne ligger värden som ändras medan programmet kör?
  \item Vilket av de två minnena behåller sitt innehåll när strömmen bryts? Varför är det viktigt
    för ett Arduino-kort?
  \item Vad heter adressen \code{0x08FF} i programmen, och vad ligger där?
\end{parts}

\exercise{Samma tal, olika skrivsätt}{Räkna för hand}
\label{c7:ex:3}
\begin{parts}
  \item Skriv tre \code{ldi}-instruktioner som var och en laddar talet 45 i \code{r16}: decimalt,
    binärt och hexadecimalt.
  \item Gör samma sak för talen 255, 16 och 100.
  \item Vilket decimalt värde laddas av \code{ldi r16, 0b01010101}, \code{ldi r17, 0x3F},
    \code{ldi r18, $80} och \code{ldi r19, 'a'}? Tecknet \code{a} har ASCII-koden \code{0x61}.
  \item En lysdiod ska lysa för varje etta i ett register. Vilket skrivsätt väljer du när du
    laddar mönstret, och varför?
\end{parts}

\exercise{Förutsäg registren}{Räkna för hand}
\label{c7:ex:4}
Fyll i tabellen med värdet i varje register \textbf{efter} varje instruktion, hexadecimalt. Alla
register är 0 när programmet startar.

\begin{asmcode}
main:
    ldi r16, 0x10
    ldi r17, 7
    mov r18, r17
    inc r18
    inc r18
    dec r16
    mov r19, r16
    clr r17
    ser r20
    dec r20
end:
    rjmp end
\end{asmcode}

\begin{filltable}{@{}p{7.6em}YYYYY@{}}
  \thead{Efter} & \thead{\code{r16}} & \thead{\code{r17}} & \thead{\code{r18}} &
    \thead{\code{r19}} & \thead{\code{r20}} \\
  \midrule
  \code{ldi r16, 0x10} & & & & & \fillrule
  \code{ldi r17, 7} & & & & & \fillrule
  \code{mov r18, r17} & & & & & \fillrule
  \code{inc r18} & & & & & \fillrule
  \code{inc r18} & & & & & \fillrule
  \code{dec r16} & & & & & \fillrule
  \code{mov r19, r16} & & & & & \fillrule
  \code{clr r17} & & & & & \fillrule
  \code{ser r20} & & & & & \fillrule
  \code{dec r20} & & & & & \\
\end{filltable}

Kontrollera sedan ditt svar genom att stega programmet i simulatorn.

\exercise{Hitta felen}{Förståelse}
\label{c7:ex:5}
Programmet nedan går inte att assemblera. Hitta alla fel, förklara vart och ett, och skriv det
rättade programmet.

\begin{asmcode}
.include "m328Pdef.inc"

.org 0x0000
    rjmp main

main
    ldi r16 5
    ldi r10, 3
    mov r17 r16
    inc r20, 1
    rjmp mian
\end{asmcode}

\exercise{Ditt första egna program}{Program}
\label{c7:ex:6}
Skriv ett program, utifrån mallen (sidan~\pageref{lab3:template}), som:
\begin{enumerate}
  \item laddar talet 12 i \code{r16};
  \item kopierar \code{r16} till \code{r17};
  \item ökar \code{r17} med 3, med \code{inc};
  \item laddar bitmönstret \code{1010 1010} i \code{r18}, skrivet binärt;
  \item kopierar \code{r18} till \code{r19};
  \item nollställer \code{r18}.
\end{enumerate}
\begin{parts}
  \item Skriv ned vilka värden \code{r16}--\code{r19} ska ha när programmet når \code{end}, både
    decimalt och hexadecimalt.
  \item Bygg och stega programmet. Stämmer värdena?
  \item Hur många klockcykler har programmet kört när det når \code{end}? Räkna först, kontrollera
    sedan med Cycle Counter.
\end{parts}

\exercise{Maskinkod för hand}{Kontroll}
\label{c7:ex:7}
\emph{Räkna ut för hand, kontrollera i listfilen, förklara.}

Använd kodningen av \code{ldi} i \secref{c7:a6}: \code{1110 KKKK dddd KKKK}, där \code{dddd} är
registrets nummer minus 16.
\begin{parts}
  \item Räkna ut maskinkoden, hexadecimalt, för:
    \begin{itemize}
      \item \code{ldi r16, 0x2A}
      \item \code{ldi r24, 0x2A}
      \item \code{ldi r31, 0xFF}
      \item \code{ldi r20, 0x03}
    \end{itemize}
    Visa de fyra fälten för åtminstone den första.
  \item Skriv de fyra instruktionerna i ett program, bygg det, och öppna listfilen (\code{.lss}) i
    projektets katalog \code{Debug}. Jämför maskinkoden där med dina svar.
  \item Stämde alla fyra? Om inte: vilket fält hade du fel på? Det vanligaste felet är att skriva
    konstanten i ett stycke i stället för att dela den i två halvor.
  \item Varför finns det ingen maskinkod för \code{ldi r15, 0x2A}?
\end{parts}

\exercise{Programmet som inte tog slut}{Förståelse}
\label{c7:ex:8}
\begin{parts}
  \item Varför slutar varje program i boken med raderna \code{end:} och \code{rjmp end}?
  \item Ta bort de två raderna ur programmet i övning~\ref{c7:ex:4}. Programmet går fortfarande att
    assemblera. Vad händer när processorn har kört den sista instruktionen?
  \item Varför börjar varje program med \code{rjmp main}, trots att \code{main} kommer direkt
    efter?
\end{parts}

\exercise{Tillåtet eller inte?}{Förståelse}
\label{c7:ex:9}
Vilka av raderna är tillåtna? Förklara varför de otillåtna inte är det, och skriv en rad som gör
det som troligen var meningen.
\begin{parts}
  \item \code{mov r5, r16}
  \item \code{ldi r5, 16}
  \item \code{mov r16, 5}
  \item \code{ldi r16, r5}
  \item \code{inc r3}
\end{parts}

\exercise{Hur lång tid tar programmet?}{Räkna för hand}
\label{c7:ex:10}
\begin{parts}
  \item Hur många klockcykler tar programmet i övning~\ref{c7:ex:4}, från start tills det når
    \code{end}? Programmet står i mallen, så det börjar med \code{rjmp main} på adress 0.
  \item Hur lång tid är det vid 16~MHz?
  \item Ungefär hur många gånger per sekund skulle processorn hinna köra programmet, om det startade
    om efter varje varv?
\end{parts}

\exercise{Maskinkod baklänges}{Räkna för hand, Fördjupning}
\label{c7:ex:11}
Tre instruktioner i en listfil har maskinkoden nedan. Alla tre är \code{ldi}-instruktioner. Vilka?
\begin{parts}
  \item \code{0xE3F5}
  \item \code{0xE0A0}
  \item \code{0xEF0F}. Vilken annan instruktion ur \secref{c7:b5} gör exakt samma sak?
\end{parts}
